\scanpage{209}
\mspage{158}

\uline{3.5}. \uline{\unsure{Application|Applications} aux sommes amalgamées dans une catégorie} $X^\wedge$ (avec $X\in\operatorname{Ob}\mathrm{Cat}$).

D'abord une caractérisation des sommes amalgamées impliquant une imm. ouverte :

\uline{Proposition 15} Soit

\diagram{\begin{tikzcd}[ampersand replacement=\&]
\& U \arrow[r,hook,"i"]\arrow[d,"{\unsure{r|p}}"'] \& X\arrow[d,"q"]\\
Q\colon \& U' \arrow[r,"i'"] \& X'
\end{tikzcd}}

un carré commutatif dans Cat. Pour que ce carré soit cocartésien\inserted{ dans Cat}, avec $i$ (donc $i'$) une imm. ouverte, il faut et il suffit que les conditions suivantes soient satisfaites :

\inserted{a) $Q$ est cartésien,}

\deleted{\unsure{a}}b) $i'$ est une immersion ouverte (donc $i$ aussi),

\deleted{b) $Q$ est cartésien}

c) Si $T'=X'\setminus U'$, et $T=X\times_{X'}T'\inserted{=X\setminus U}$, $T\longrightarrow T'$ est un iso.

d) Pour tout $t\in T$, si $t'$ est son image dans $T'$, le foncteur induit par $p$

\[
p_{/t}\colon
\underset{\substack{\shortparallel\\ V_t}}{U_{/t}}
\longrightarrow
\underset{\substack{\shortparallel\\ V'_{t'}}}{U'_{/t'}}
\]

\deleted{\unsure{satis}} \deleted{\unsure{est à}} satisfait : $(p_{/t})_!(e_{U/t})\simeq e_{U'/t'}$, ou ce qui revient au $=$, les ``fibres colocales'' de $p_{/t}$, $v'\backslash V_t$ avec $v'$ dans $V'_{t'}$, sont $0$-connexes.
% archive sheet number: 209
